\subsection{Evaluación Exploratoria de Pruebas \textit{End-to-End}}\label{subsec:evaluacion-e2e}

Durante una fase exploratoria se evaluó incorporar pruebas \textit{end-to-end} (E2E) con Playwright sobre los emuladores locales de Firestore y Authentication, para valorar la viabilidad de automatizar los recorridos de interfaz verificados hasta entonces de forma manual. La exploración estructuró los cimientos de una \textit{suite} —configuración contra los emuladores, patrón \textit{Page Object Model}, \textit{fixtures} para preparar el estado de Firebase y casos iniciales de inicio de sesión y registro—, pero su adopción se mantuvo a un nivel mínimo y no se consolidó como parte de la \textit{suite} oficial del proyecto.

La decisión respondió a un criterio de costo-beneficio: las pruebas E2E de interfaz son las más costosas de mantener de la pirámide —sensibles a cambios de interfaz, propensas a fallos intermitentes y de mantenimiento sostenido—, difíciles de absorber por un equipo reducido sin desplazar esfuerzo de las funcionalidades prioritarias. Por ello se concentró la automatización en la capa de mayor riesgo y menor volatilidad —las reglas de seguridad— y se sostuvo la validación de extremo a extremo con el proceso manual descrito. Documentar esta evaluación responde a un criterio de transparencia: dejó una base reutilizable para una futura consolidación de la \textit{suite} E2E, posibilidad que se retoma en el trabajo futuro.
